% Chronological networks over an interval-native date/time library.
%
% Format: OASIcs (oasics-v2021), the TIME-symposium template, reused
% from ../time-2026. The class file and its logos (cc-by.pdf,
% oasics-logo-bw.pdf, orcid.pdf) sit alongside this file.
%
% Build:
%   pdflatex chrononetworks && pdflatex chrononetworks
%
% Uses an inline thebibliography so no .bib is required at draft stage.

\documentclass[a4paper,USenglish,cleveref,autoref,thm-restate]{oasics-v2021}

\usepackage{amssymb}
\usepackage{booktabs}

% The OASIcs class defines \doi for BibTeX; provide it for the inline
% thebibliography below (mirrors the plainurl style).
\providecommand{\doi}[1]{\href{https://doi.org/#1}{doi:\texttt{#1}}}

\title{Chronological Networks as a Library Primitive: \texorpdfstring{\\}{ }%
  ChronoLog over an Interval-Native Date/Time Stack}
\titlerunning{Chronological Networks as a Library Primitive}

\author{Kip Cole}{Independent}{kipcole9@gmail.com}{}{}
\authorrunning{K.~Cole}
\Copyright{Kip Cole}

\ccsdesc[500]{Applied computing~Arts and humanities}
\ccsdesc[300]{Theory of computation~Modal and temporal logics}
\ccsdesc[300]{Information systems~Temporal data}

\keywords{Chronological networks, Allen's interval algebra, simple
  temporal problems, archaeological dating, ISO 8601, EDTF,
  interval-based time}

\relatedversion{Source repository: \url{https://github.com/kipcole9/tempo}}

\begin{document}

\maketitle

\begin{abstract}
  Levy, Geeraerts, Pluquet, Piasetzky and Fantalkin (2020) formalised
  archaeological chronology as a \emph{chronological network} --- a set
  of time-periods related by synchronisms, asynchronisms, and duration
  bounds --- and showed that consistency checking and the computation
  of tightest possible date ranges reduce to a Simple Temporal Problem
  (STP), solvable in polynomial time. Their tool, ChronoLog, invents an
  interval/relation vocabulary from scratch. We observe that this
  vocabulary is exactly the one a general, interval-native date/time
  library already provides: half-open intervals, Allen's thirteen
  relations, partial and reduced-precision dates, calendar-aware
  durations, and uncertainty qualifiers. We implement the ChronoLog
  model as a small layer (\texttt{Tempo.Network}) over the Tempo Elixir
  library, mapping each ChronoLog relation to its Allen counterpart plus
  metric (delay) extensions, normalising a network to an STP, and
  solving consistency and tightening by Floyd--Warshall. The layer
  reproduces the paper's published bounds exactly and inherits, for
  free, ISO~8601 / EDTF interchange, multiple calendars, BCE arithmetic,
  and §8 uncertainty qualifiers. The contribution is less a new
  algorithm than an argument: chronological-network modelling is a
  natural \emph{library primitive} when the library already treats time
  as intervals.
\end{abstract}

\section{Introduction}

Archaeological and historical chronology is rarely a list of fixed
dates. It is a web of relationships: a king reigns before another, a
stratum is built during a reign, a ceramic style overlaps its
successor, an inscription supplies a \emph{terminus post quem}. Levy et
al.~\cite{levy2020} call such a web a \emph{chronological network} and
give it a formal model with three object types --- time-periods,
sequences, and chronological relations --- and two operations:
\emph{consistency checking} (does any assignment of dates satisfy every
constraint?) and \emph{tightening} (what is the narrowest start, end,
and duration each period can take?). Their key technical result is that
the model is a Simple Temporal Problem~\cite{dechter1991}: every
relation, duration, and absolute date reduces to atomic constraints of
the form $b_1 - b_2 \leq k$ over boundary variables, so consistency is
the absence of a negative cycle and tightening is all-pairs shortest
paths. This is polynomial-time, in contrast to exhaustive-search tools
such as Groundhog~\cite{falk2020}.

ChronoLog~\cite{chronolog} implements this model as a bespoke
application, and in doing so re-creates a small interval algebra: a
period has a start, an end, and a duration, each possibly unknown,
bounded, or ranged; relations express containment, overlap, synchronised
boundaries, and metric delays. Our observation is that this is precisely
the vocabulary that an \emph{interval-native} general-purpose date/time
library already provides. We demonstrate this with Tempo, an Elixir
library in which every temporal value is a bounded half-open interval at
a declared resolution, comparison is Allen's thirteen
relations~\cite{allen1983}, and partial ISO~8601 / EDTF values
(\texttt{2022}, \texttt{2022-06}, \texttt{$\sim$720}) are first-class.

We contribute \texttt{Tempo.Network}, a layer that (i) names the
ChronoLog relation vocabulary and maps it to Allen's algebra plus metric
extensions (§\ref{sec:mapping}); (ii) normalises a network to an STP
against a single origin boundary (§\ref{sec:solver}); and (iii) solves
consistency and tightening by Floyd--Warshall, with shortest-path
\emph{traces} that explain each derived bound (§\ref{sec:solver}). The
layer reproduces the paper's ChronoLand and 26th-dynasty results exactly
(§\ref{sec:examples}) and inherits ISO~8601/EDTF interchange, calendars,
and qualifiers from the host library (§\ref{sec:beyond}).

\section{The ChronoLog model}
\label{sec:chronolog}

A \emph{time-period} is a continuous interval --- a reign, an era, the
life of a stratum --- characterised by a start date, an end date, and a
duration, each of which may be unknown, known, lower-bounded,
upper-bounded, or known within a range $[lo, hi]$. Dates and durations
are modelled \emph{independently}: a period is at most six numbers
(earliest/latest start, earliest/latest end, minimum/maximum duration),
reconciled by the relation $\textit{duration} = \textit{end} -
\textit{start}$.

A \emph{sequence} is a list of consecutive time-periods with no gaps:
$\textit{end}(p_i) = \textit{start}(p_{i+1})$. A gap is modelled by
inserting an explicit gap period.

A \emph{chronological relation} constrains two periods. The vocabulary
(the paper's Tables~1--5) covers contemporaneity ($A \sim B$: non-empty
overlap), inclusion ($A \subseteq B$, $A \supseteq B$), overlap ($A \leq
B$, $A \geq B$), start/end-period synchronisms (``$A$ starts during
$B$'', ``$A$ includes the start of $B$'', and the end-side duals),
synchronised boundaries (synchronous start/end, equality, immediate
precedence/following), asynchronisms ($A \ll B$, $A \gg B$), ordered
boundaries (strict orderings of individual start/end dates), and
\emph{delay synchronisms} --- metric constraints stating that one
boundary is an exact, minimum, or maximum number of years before or
after another.

\paragraph{The STP reduction.} Every relation, duration, sequence link,
and absolute date normalises to a conjunction of atomic constraints of
one shape, $b_1 - b_2 \leq k$, where $b_1, b_2$ are boundary variables
and $k$ is an integer. Equalities become two inequalities; a strict $<$
on an integer time-scale becomes $\leq k-1$; absolute dates are
constraints against a single origin boundary $z_0$. The constraints form
a directed weighted graph with one node per boundary plus $z_0$. The
network is \emph{consistent} iff the graph has no negative cycle, and the
tightest bound on $b_1 - b_2$ is the shortest-path weight $b_1 \to b_2$;
tightening is therefore all-pairs shortest paths
(Floyd--Warshall)~\cite{dechter1991,floyd1962}.

\section{Tempo's interval model}
\label{sec:tempo}

Tempo represents every temporal value as a bounded, half-open interval
$[\textit{from}, \textit{to})$ at a declared resolution. A partial
ISO~8601 value is an interval of the appropriate width: \texttt{2026} is
$[2026\text{-}01\text{-}01, 2027\text{-}01\text{-}01)$, \texttt{2026-06}
spans one month, \texttt{2026-06-15} one day. The half-open convention
makes adjacent spans concatenate without gap or overlap, and makes
comparison total under Allen's thirteen relations, exposed as predicates
(\texttt{before?}, \texttt{overlaps?}, \texttt{within?},
\texttt{adjacent?}). Durations are calendar-aware
(\texttt{P20Y}, \texttt{P2Y6M}), and ISO~8601-2 / EDTF uncertainty
qualifiers ($?$ uncertain, $\sim$ approximate, $\%$ both) are parsed and
carried per component.

These are exactly the primitives ChronoLog needs. A time-period's
boundaries are interval endpoints; a sequence is the half-open
concatenation $[a,b)$ then $[b,c)$; the relations are Allen relations
plus metric offsets; absolute and floating chronologies are anchored and
unanchored values; uncertainty is the qualifier vocabulary. The one
missing piece is the constraint solver, which is small and
well-understood. \texttt{Tempo.Network} supplies it.

\section{Mapping the relation vocabulary}
\label{sec:mapping}

Each ChronoLog relation is stored as data and translated, by
\texttt{to\_atomic/1}, to its atomic constraints; Table~\ref{tab:map}
gives the correspondence to Allen's algebra and the atomic form. Where a
single Allen relation exists the mapping is one-to-one (and invertible,
so a solved network can be \emph{queried} with Allen predicates and an
Allen relation between two solved periods can be \emph{named}). The
period synchronisms (``starts during'' and kin) are looser than any
single Allen relation --- they are disjunctions --- and the metric
delays carry an integer Allen cannot express.

\begin{table}[t]
  \caption{Selected ChronoLog relations, their Allen counterpart, and
    their atomic constraints. $s,e$ abbreviate start, end.}
  \label{tab:map}
  \centering
  \begin{tabular}{lll}
    \toprule
    ChronoLog relation & Allen & Atomic constraints \\
    \midrule
    $A \ll B$ (before)            & precedes   & $e(A) - s(B) \leq -1$ \\
    $A \sqsupset B$ (immed.\ prec.) & meets    & $e(A) = s(B)$ \\
    $A \subseteq B$ (included in) & during     & $s(B) \leq s(A),\ e(A) \leq e(B)$ \\
    $A \supseteq B$ (includes)    & contains   & $s(A) \leq s(B),\ e(B) \leq e(A)$ \\
    $A \leq B$ (overlaps)         & overlaps   & $s(A) \leq s(B) \leq e(A) \leq e(B)$ \\
    $A \sim B$ (contemporary)     & ---        & $s(B) \leq e(A),\ s(A) \leq e(B)$ \\
    $A$ starts during $B$         & ---        & $s(B) \leq s(A) \leq e(B)$ \\
    $A$ ends during $B$           & ---        & $s(B) \leq e(A) \leq e(B)$ \\
    $A \top B$ (synchronous start) & ---       & $s(A) = s(B)$ \\
    delay: $s(A)$ at least $X$ before $s(B)$ & --- & $s(A) - s(B) \leq -X$ \\
    \bottomrule
  \end{tabular}
\end{table}

\section{Normalisation and solver}
\label{sec:solver}

\texttt{Tempo.Network.Normalize} reduces a network to
$\{\textit{nodes}, \textit{edges}, \textit{unit}\}$. Each period
contributes \texttt{start} and \texttt{end} nodes; a shared origin
\texttt{:origin} ($z_0$) anchors absolute dates. The network's
\emph{finest unit} --- the finest resolution among its dates --- fixes
the integer axis: a date becomes its count of that unit relative to
$z_0$, a duration its count of that unit. Conversion is exact for the
uniform-unit (typically year-grained) chronologies that are the
discipline's norm; a duration expressed across units uses nominal
calendar lengths. Each edge is tagged with the constraint that produced
it, for tracing.

\texttt{Tempo.Network.Solver} runs Floyd--Warshall once, in $O(n^3)$ on
the boundary count --- interactive for the hundreds of periods these
chronologies contain. \texttt{consistent?/1} reports the absence of a
negative cycle; \texttt{tighten/1} reads the tightest start, end, and
duration of each period off the shortest-path matrix and maps them back
to Tempo values. \texttt{trace/3} reconstructs the shortest path
underlying a chosen bound and renders it as prose --- the paper's notion
of an explanatory trace, where every step names the input it rests on.

\section{Worked examples}
\label{sec:examples}

\paragraph{ChronoLand.} The paper's toy kingdom: kings \texttt{K1} then
\texttt{K2} reign in succession, both between 1200 and 1300~CE; \texttt{K1}
reigns at most 10 years, \texttt{K2} at least 35. Strata \texttt{S1}
(built by \texttt{K1}) and \texttt{S2} (destroyed under \texttt{K2})
follow one another, each lasting 20--100 years.

\begin{verbatim}
network =
  Network.new()
  |> Network.add_period(:k1, start: {:not_before, ~o"1200Y"}, duration: {:at_most, ~o"P10Y"})
  |> Network.add_period(:k2, end: {:not_after, ~o"1300Y"}, duration: {:at_least, ~o"P35Y"})
  |> Network.add_period(:s1, duration: {~o"P20Y", ~o"P100Y"})
  |> Network.add_period(:s2, duration: {~o"P20Y", ~o"P100Y"})
  |> Network.add_sequence([:k1, :k2])
  |> Network.add_sequence([:s1, :s2])
  |> Network.add_relation(:starts_during, :s1, :k1)
  |> Network.add_relation(:ends_during, :s2, :k2)
\end{verbatim}

Tightening reproduces the paper's Fig.~6b exactly: \texttt{K2} tightens
to start $\in [1200, 1265]$, end $\in [1240, 1300]$, duration $\in [35,
100]$; the strata, given no dates as input, acquire them, and each
stratum's duration tightens from $[20,100]$ to $[20,80]$ because the
dynasty is too short to hold a longer one. The trace for the earliest
end of \texttt{K2} reproduces the paper's Fig.~6c reasoning:
\texttt{K1} starts no earlier than 1200; \texttt{S1} starts during
\texttt{K1}; \texttt{S1} lasts at least 20 years, so ends after 1220;
\texttt{S1} meets \texttt{S2}; \texttt{S2} lasts at least 20 years, so
ends after 1240; \texttt{S2} ends during \texttt{K2}, so \texttt{K2}
ends after 1240. Capping \texttt{K2} at 25 rather than 35 years makes
the network inconsistent (Fig.~7), detected as a negative cycle.

\paragraph{The 26th dynasty.} Six reigns in succession (after Kitchen
2000). Given only Psammetichus~I's accession (664~BCE) and each reign's
length, tightening derives every reign's absolute dates --- exercising
BCE (negative) years and a six-period sequence. Amasis~II's accession is
recovered as 570~BCE, forced by the anchor and the five preceding
reigns, with no date stated for it directly.

\section{What the library adds, and its limits}
\label{sec:beyond}

Because boundaries are ordinary Tempo values, a network is interoperable
by construction: each period round-trips through ISO~8601 / EDTF for
exchange with other tooling, BCE and proleptic dates are handled
uniformly, and non-Gregorian calendars are available through the
underlying calendrical stack. Uncertainty qualifiers are carried as
metadata and, by design, do \emph{not} silently widen a bound: a vague
``circa'' never performs arithmetic unless the modeller restates it as a
range. The solver shares interval-arithmetic primitives with the
library's planned set-operations milestone (union, intersection,
coalescing of interval sets).

\paragraph{Limitations.} Following the source model, relations compose
by \emph{conjunction}; a genuine \emph{or}-relation (``$A$ ends before
$B$ \emph{or} $A$ starts after $B$'') is outside the STP and is not
supported. There is no fuzzy or Bayesian layer; uncertainty is bounds,
not probability. Cross-unit duration conversion (e.g.\ a year-grained
duration on a day axis) uses nominal calendar lengths and is therefore
approximate, though the uniform-unit chronologies that dominate practice
are exact.

\section{Related work}

Allen's interval algebra~\cite{allen1983} underlies the relation
vocabulary; Holst~\cite{holst2004} applied temporal logic to archaeology.
The direct antecedent is ChronoLog~\cite{levy2020,chronolog}; the STP
machinery is that of Dechter, Meiri and Pearl~\cite{dechter1991} via
Floyd~\cite{floyd1962}. Falk's Groundhog~\cite{falk2020} addresses the
same problem by exhaustive search. Tempo's contribution is orthogonal to
all of these: it shows that, given an interval-native host library, the
network model and its solver are a compact, reusable primitive rather
than a standalone application.

\section{Conclusion}

Chronological-network modelling does not require a bespoke tool. When the
host library already represents time as half-open intervals with Allen's
algebra, partial dates, and calendar-aware durations, the ChronoLog
model is a thin layer: a relation-to-atomic translation, an STP
normalisation, and a Floyd--Warshall solver. The result reproduces the
reference results exactly while inheriting ISO~8601/EDTF interchange,
calendars, and uncertainty qualifiers --- and makes consistency checking
and tightening available wherever the library is.

%
% References
%
\begin{thebibliography}{10}

\bibitem{allen1983}
James F. Allen.
\newblock Maintaining knowledge about temporal intervals.
\newblock \emph{Communications of the ACM}, 26(11):832--843, 1983.
\newblock \doi{10.1145/182.358434}.

\bibitem{dechter1991}
Rina Dechter, Itay Meiri, and Judea Pearl.
\newblock Temporal constraint networks.
\newblock \emph{Artificial Intelligence}, 49(1--3):61--95, 1991.
\newblock \doi{10.1016/0004-3702(91)90006-6}.

\bibitem{falk2020}
Dan Falk.
\newblock Groundhog: A chronological modelling tool.
\newblock \url{http://www.groundhogchronology.com/}, 2020.

\bibitem{floyd1962}
Robert W. Floyd.
\newblock Algorithm 97: Shortest path.
\newblock \emph{Communications of the ACM}, 5(6):345, 1962.
\newblock \doi{10.1145/367766.368168}.

\bibitem{holst2004}
Mads K. Holst.
\newblock Complicated relations and blind dating: Formal analysis of
  relative chronological structures.
\newblock In \emph{Tools for Constructing Chronologies}, pages 129--147.
  Springer, 2004.

\bibitem{levy2020}
Eythan Levy, Gilles Geeraerts, Fr\'ed\'eric Pluquet, Eli Piasetzky, and
  Alexander Fantalkin.
\newblock Chronological networks in archaeology: A formalised scheme.
\newblock \emph{Journal of Archaeological Science}, 124:105225, 2020.
\newblock \doi{10.1016/j.jas.2020.105225}.

\bibitem{chronolog}
ChronoLog.
\newblock \url{http://chrono.ulb.be}.

\end{thebibliography}

\end{document}
